% Appendix B -- the 60-minute master benchmark and the controls that say what
% the classical baselines' forewarning columns actually measure.
\section{Extended Benchmark Results (\texorpdfstring{$\PH = 60$}{PH = 60}) and Baseline Controls}
\label{app:benchmarks}

\paragraph*{How to read the classification columns of Tables~\ref{tab:sota_ph30}
and~\ref{tab:sota_ph60}.} Persistence, Ridge-AR and ARIMA have no classification
head. Their PR-AUC, Se@FPR and FAPD entries come from
$\sigma((70 - \hat{y})/\tau)$ applied to a forecast of $g(t{+}\mathrm{PH})$,
while the label is $\min g$ over $(t, t{+}\mathrm{PH}]$. Those three entries are
therefore lower bounds on what the method achieves when fitted to the labelled
quantity, and Table~\ref{tab:ablation_f_nadir_target_ph30} refits Ridge-AR to the
nadir to quantify the gap. LightGBM and both neural arms fit a classifier on the
nadir label directly and are unaffected. This caveat applies to every benchmark
table in the paper; the captions state it in one clause and point here.

\subsection{Sixty-Minute Horizon: Master Benchmark}
\label{sec:appendix_benchmarks}

% These two generated tables have sixteen columns. Set as ordinary table*
% floats they are shrunk to about 5 pt; set landscape (rotating's
% sidewaystable*, inside a group so no other float changes) they print near
% their natural size. The generated files are not edited.
\begingroup
\renewenvironment{table*}[1][]{\begin{sidewaystable*}[p]}{\end{sidewaystable*}}
\begin{table*}[t]
\centering
\small
\setlength{\tabcolsep}{4pt}
\caption{OhioT1DM official test split, PH = 30 min (12 subjects, $n = 28335$ windows). Coverage and width are split-conformal at $\alpha = 0.1$, calibrated on the disjoint calibration split; hypo coverage is the Mondrian $<70$ mg/dL stratum. Persistence, Ridge-AR and ARIMA have no classification head, so their PR-AUC, Se@FPR and FAPD entries are \emph{lower bounds} (Appendix~\ref{app:benchmarks}).}
\label{tab:sota_ph30}
\resizebox{\ifdim\width>\linewidth\linewidth\else\width\fi}{!}{%
\begin{tabular}{lrrrrrrrrrrrrrrr}
\toprule
Model & RMSE & MAE & MARD (\%) & Lag (min) & A+B (\%) & D+E ($n$) & Parkes A+B (\%) & PR-AUC & Se@5\% & Se@10\% & FAPD@5\% & Onset lead (min) & Cov. & Hypo cov. & Width \\
\midrule
Persistence & 23.29 & 16.79 & 11.25 & 30.0 & 99.02 & 272 & 99.66 & 0.832 & 0.933 & 0.970 & 1.16 & 23.7 & 0.913 & \textbf{0.934} & 80.0 \\
Ridge-AR & 20.05 & 14.29 & 9.76 & 25.0 & 98.74 & 351 & 99.76 & 0.829 & 0.964 & 0.990 & 2.99 & 26.4 & 0.916 & 0.933 & 67.7 \\
ARIMA(4,1,0) & 21.09 & 14.89 & 9.94 & 25.0 & 99.31 & 187 & 99.79 & 0.856 & 0.957 & 0.985 & 2.05 & 25.8 & 0.914 & 0.933 & 72.5 \\
LightGBM & \textbf{18.92} & 13.07 & 8.74 & \textbf{20.0} & 99.04 & 270 & 99.81 & 0.841 & \textbf{0.968} & 0.993 & 2.90 & \textbf{27.1} & 0.920 & 0.922 & 66.6 \\
Dilated LSTM (CGM) & 19.52 & 13.39 & 8.80 & 20.0 & \textbf{99.38} & \textbf{175} & \textbf{99.86} & 0.858 & 0.962 & \textbf{0.994} & 2.88 & 26.2 & 0.918 & 0.933 & 68.5 \\
CausalCGM-MT (CGM) & 19.16 & 13.33 & 8.89 & 20.0 & 99.18 & 231 & 99.79 & 0.868 & 0.954 & 0.990 & 2.57 & 25.6 & 0.922 & 0.926 & 68.7 \\
CausalCGM-MT (fusion) & 18.97 & \textbf{13.01} & \textbf{8.52} & 20.0 & 99.34 & 184 & 99.84 & \textbf{0.870} & 0.957 & 0.991 & 2.80 & 26.5 & 0.926 & 0.887 & 68.8 \\
CausalCGM-MT + LoRA ($r = 4$) & 21.81 & 15.82 & 10.52 & 25.0 & 98.86 & 323 & 99.67 & 0.649 & 0.783 & 0.898 & 1.80 & 21.9 & 0.921 & 0.907 & 78.3 \\
\bottomrule
\end{tabular}}
\end{table*}

\begin{table*}[t]
\centering
\small
\setlength{\tabcolsep}{4pt}
\caption{OhioT1DM official test split, PH = 60 min (12 subjects, $n = 28335$ windows). Coverage and width are split-conformal at $\alpha = 0.1$, calibrated on the disjoint calibration split; hypo coverage is the Mondrian $<70$ mg/dL stratum. Persistence, Ridge-AR and ARIMA have no classification head, so their PR-AUC, Se@FPR and FAPD entries are \emph{lower bounds} (Appendix~\ref{app:benchmarks}).}
\label{tab:sota_ph60}
\resizebox{\ifdim\width>\linewidth\linewidth\else\width\fi}{!}{%
\begin{tabular}{lrrrrrrrrrrrrrrr}
\toprule
Model & RMSE & MAE & MARD (\%) & Lag (min) & A+B (\%) & D+E ($n$) & Parkes A+B (\%) & PR-AUC & Se@5\% & Se@10\% & FAPD@5\% & Onset lead (min) & Cov. & Hypo cov. & Width \\
\midrule
Persistence & 38.16 & 28.12 & 19.03 & 60.0 & 96.88 & 700 & 97.84 & 0.724 & 0.789 & 0.866 & 1.25 & 35.8 & 0.909 & 0.928 & 129.9 \\
Ridge-AR & 33.82 & 25.32 & 17.63 & 50.0 & 96.27 & 987 & 98.01 & 0.702 & 0.864 & 0.920 & 3.72 & 45.1 & 0.905 & 0.936 & 112.3 \\
ARIMA(4,1,0) & 36.55 & 26.68 & 17.95 & 55.0 & 97.25 & \textbf{611} & 98.18 & 0.752 & 0.830 & 0.901 & 2.34 & 40.8 & 0.911 & 0.931 & 125.6 \\
LightGBM & 32.12 & 23.13 & 15.53 & 45.0 & 96.93 & 820 & 98.50 & 0.755 & \textbf{0.867} & 0.926 & 3.76 & \textbf{46.0} & 0.910 & 0.939 & 110.5 \\
Dilated LSTM (CGM) & 33.50 & 23.89 & 15.54 & 45.0 & 97.50 & 662 & \textbf{98.81} & 0.766 & 0.858 & 0.920 & 3.28 & 44.5 & 0.911 & \textbf{0.952} & 117.2 \\
CausalCGM-MT (CGM) & 33.07 & 23.70 & 15.80 & 45.0 & 97.21 & 711 & 98.64 & 0.775 & 0.844 & 0.910 & 2.99 & 42.1 & 0.918 & 0.942 & 118.5 \\
CausalCGM-MT (fusion) & \textbf{31.76} & \textbf{22.64} & \textbf{14.83} & \textbf{40.0} & \textbf{97.55} & 676 & 98.78 & \textbf{0.783} & 0.858 & \textbf{0.928} & 3.00 & 44.7 & 0.925 & 0.947 & 117.4 \\
CausalCGM-MT + LoRA ($r = 4$) & 37.59 & 28.01 & 18.72 & 45.0 & 96.96 & 824 & 97.67 & 0.548 & 0.589 & 0.731 & 1.63 & 32.9 & 0.927 & 0.938 & 137.4 \\
\bottomrule
\end{tabular}}
\end{table*}

\endgroup

Tables~\ref{tab:sota_ph30} and~\ref{tab:sota_ph60} give every metric for one
seed ($1337$) at $\PH = 30$ and $60$~min on the official test split. They exist
for completeness: a bold entry marks the best value in that single run and is
not evidence of a difference. Comparisons between arms should be read from the
multi-seed tables below. Persistence, Ridge-AR and ARIMA have no classification
head, so their forewarning columns are lower bounds, which
Appendix~\ref{sec:ablation_f} quantifies.

% A full-page rotated float otherwise waits at the head of the two-column
% float queue and holds every later table back to the end of the paper; flush
% the two landscape tables here, before the next subsection.
\clearpage
\subsection{Multi-Seed Results and Bootstrap Intervals}
\label{sec:appendix_multiseed}

\begin{table*}[t]
\centering
\small
\setlength{\tabcolsep}{4pt}
\caption{Fusion minus each arm at $\PH = 30$~min, seed-mean difference with a 95\% subject-cluster bootstrap CI (12 subjects resampled, 1000 draws). A CI spanning 0 is not a resolved difference.}
\label{tab:multiseed_diffs_ph30}
\resizebox{\ifdim\width>\linewidth\linewidth\else\width\fi}{!}{%
\begin{tabular}{lrrrrr}
\toprule
vs.\ arm & $\Delta$RMSE & $\Delta$PR-AUC & $\Delta$PR-AUC (onset) & $\Delta$FAPD & $\Delta$(D+E) \\
\midrule
Persistence & -4.53 [-5.36, -3.70] & +0.032 [-0.008, +0.069] & +0.104 [-0.006, +0.203] & +1.31 [+0.90, +1.73] & -75 [-123, -28] \\
Ridge-AR & -1.29 [-1.62, -0.88] & +0.034 [+0.014, +0.063] & +0.021 [-0.036, +0.082] & -0.52 [-0.99, -0.08] & -154 [-229, -80] \\
ARIMA(4,1,0) & -2.33 [-2.67, -1.95] & +0.008 [-0.022, +0.041] & +0.032 [-0.062, +0.133] & +0.42 [+0.12, +0.73] & +10 [-27, +47] \\
LightGBM & -0.15 [-0.48, +0.14] & +0.022 [-0.012, +0.053] & -0.023 [-0.098, +0.052] & -0.43 [-0.88, -0.15] & -73 [-125, -24] \\
Dilated LSTM (CGM) & -0.59 [-0.96, -0.24] & +0.000 [-0.019, +0.014] & +0.002 [-0.061, +0.053] & -0.14 [-0.53, +0.11] & +24 [-3, +59] \\
CausalCGM-MT (CGM) & -0.52 [-0.82, -0.23] & +0.001 [-0.013, +0.009] & +0.023 [-0.017, +0.052] & -0.17 [-0.54, +0.14] & +2 [-21, +30] \\
CausalCGM-MT + LoRA (r=4) & -2.85 [-3.51, -2.34] & +0.193 [+0.136, +0.224] & +0.190 [+0.128, +0.235] & +0.69 [-0.11, +1.45] & -145 [-225, -69] \\
\bottomrule
\end{tabular}}
\end{table*}

\begin{table*}[t]
\centering
\small
\setlength{\tabcolsep}{4pt}
\caption{OhioT1DM test split, $\PH = 60$~min: mean $\pm$ sd over 5 training seeds. Alarm thresholds are fitted on the calibration split and applied unchanged to test. PR-AUC (onset) excludes anchors that are below 70~mg/dL now or in the preceding 30~min. Onset lead is minutes from the first valid alert to the first reading below 70. FAPD@5\% counts false alarm episodes under the default alarm policy: an episode lasts while the score stays above threshold, with no merging across gaps and a minimum duration of one reading (5~min); Appendix~\ref{sec:appendix_latching} varies both. Classical arms and LightGBM are deterministic (sd 0).}
\label{tab:sota_multiseed_ph60}
\resizebox{\ifdim\width>\linewidth\linewidth\else\width\fi}{!}{%
\begin{tabular}{lrrrrrrr}
\toprule
Arm & RMSE & PR-AUC & PR-AUC (onset) & Se@5\% (cal.) & FAPD@5\% & Onset lead & D+E \\
\midrule
Persistence & 38.16 $\pm$ 0.00 & 0.724 $\pm$ 0.000 & 0.312 $\pm$ 0.000 & 0.789 $\pm$ 0.000 & 1.25 $\pm$ 0.00 & 35.8 $\pm$ 0.0 & 700 $\pm$ 0 \\
Ridge-AR & 33.82 $\pm$ 0.00 & 0.702 $\pm$ 0.000 & 0.377 $\pm$ 0.000 & 0.864 $\pm$ 0.000 & 3.72 $\pm$ 0.00 & 45.1 $\pm$ 0.0 & 987 $\pm$ 0 \\
ARIMA(4,1,0) & 36.55 $\pm$ 0.00 & 0.752 $\pm$ 0.000 & 0.364 $\pm$ 0.000 & 0.830 $\pm$ 0.000 & 2.34 $\pm$ 0.00 & 40.8 $\pm$ 0.0 & 611 $\pm$ 0 \\
LightGBM & 32.12 $\pm$ 0.00 & 0.755 $\pm$ 0.000 & 0.431 $\pm$ 0.000 & 0.867 $\pm$ 0.000 & 3.76 $\pm$ 0.00 & 46.0 $\pm$ 0.0 & 820 $\pm$ 0 \\
Dilated LSTM (CGM) & 33.30 $\pm$ 0.31 & 0.773 $\pm$ 0.009 & 0.408 $\pm$ 0.022 & 0.859 $\pm$ 0.001 & 3.24 $\pm$ 0.27 & 44.7 $\pm$ 0.7 & 643 $\pm$ 21 \\
CausalCGM-MT (CGM) & 33.32 $\pm$ 0.23 & 0.774 $\pm$ 0.003 & 0.401 $\pm$ 0.016 & 0.846 $\pm$ 0.013 & 3.10 $\pm$ 0.14 & 43.1 $\pm$ 1.3 & 690 $\pm$ 49 \\
CausalCGM-MT (fusion) & 31.85 $\pm$ 0.37 & 0.774 $\pm$ 0.009 & 0.413 $\pm$ 0.022 & 0.847 $\pm$ 0.009 & 2.84 $\pm$ 0.32 & 43.9 $\pm$ 0.7 & 617 $\pm$ 53 \\
CausalCGM-MT + LoRA (r=4) & 37.62 $\pm$ 0.95 & 0.564 $\pm$ 0.065 & 0.237 $\pm$ 0.035 & 0.622 $\pm$ 0.069 & 1.78 $\pm$ 0.21 & 34.1 $\pm$ 1.8 & 859 $\pm$ 31 \\
\bottomrule
\end{tabular}}
\end{table*}

\begin{table*}[t]
\centering
\small
\setlength{\tabcolsep}{4pt}
\caption{Fusion minus each arm at $\PH = 60$~min, seed-mean difference with a 95\% subject-cluster bootstrap CI (12 subjects resampled, 1000 draws). A CI spanning 0 is not a resolved difference.}
\label{tab:multiseed_diffs_ph60}
\resizebox{\ifdim\width>\linewidth\linewidth\else\width\fi}{!}{%
\begin{tabular}{lrrrrr}
\toprule
vs.\ arm & $\Delta$RMSE & $\Delta$PR-AUC & $\Delta$PR-AUC (onset) & $\Delta$FAPD & $\Delta$(D+E) \\
\midrule
Persistence & -6.31 [-7.77, -5.04] & +0.050 [-0.003, +0.095] & +0.101 [-0.014, +0.200] & +1.59 [+1.10, +2.06] & -83 [-171, +11] \\
Ridge-AR & -1.97 [-2.92, -1.01] & +0.071 [+0.036, +0.110] & +0.036 [-0.024, +0.099] & -0.88 [-1.61, -0.33] & -370 [-649, -170] \\
ARIMA(4,1,0) & -4.70 [-5.78, -3.66] & +0.022 [-0.020, +0.063] & +0.049 [-0.041, +0.139] & +0.51 [+0.01, +0.91] & +6 [-66, +98] \\
LightGBM & -0.27 [-1.00, +0.50] & +0.018 [-0.015, +0.047] & -0.019 [-0.080, +0.040] & -0.92 [-1.66, -0.40] & -203 [-315, -106] \\
Dilated LSTM (CGM) & -1.45 [-2.32, -0.59] & +0.000 [-0.022, +0.016] & +0.004 [-0.035, +0.032] & -0.40 [-0.91, -0.01] & -26 [-73, +31] \\
CausalCGM-MT (CGM) & -1.47 [-2.20, -0.77] & +0.000 [-0.023, +0.016] & +0.012 [-0.040, +0.047] & -0.25 [-0.71, +0.15] & -73 [-120, -23] \\
CausalCGM-MT + LoRA (r=4) & -5.77 [-7.14, -4.57] & +0.210 [+0.142, +0.268] & +0.175 [+0.132, +0.220] & +1.06 [+0.18, +2.01] & -242 [-375, -129] \\
\bottomrule
\end{tabular}}
\end{table*}


Table~\ref{tab:sota_multiseed_ph30} (main text) and
Table~\ref{tab:sota_multiseed_ph60} report means over $\NumMSSeeds$ seeds;
classical arms and LightGBM are deterministic. CIs resample the twelve subjects,
so they reflect between-subject variation, and twelve clusters make them
approximate. Three readings follow. Fusion improves RMSE over the CGM-only
network at both horizons, with CIs excluding zero. Fusion and LightGBM are not
separated on RMSE or PR-AUC at either horizon. The three neural arms are not
separated on PR-AUC, so no arm leads on forewarning.

\subsection{Sensitivity of False-Alarm Burden to the Alarm Policy}
\label{sec:appendix_latching}

\begin{table*}[t]
\centering
\small
\setlength{\tabcolsep}{4pt}
\caption{False alarms per day / onset detection rate at $\PH = 30$~min under alternative alarm policies, mean over 5 seeds. G: latching gap (min) within which a new above-threshold run does not sound again; D: minimum duration (min) before an alarm sounds. G=0, D=5 is the policy used elsewhere in the paper. Thresholds are each seed's calibration-split 5\%-FPR thresholds, not refitted per policy.}
\label{tab:fapd_latching}
\resizebox{\ifdim\width>\linewidth\linewidth\else\width\fi}{!}{%
\begin{tabular}{lrrrrrrr}
\toprule
Policy (G, D) & Persist. & Ridge & ARIMA & LGBM & LSTM & TCN-CGM & Fusion \\
\midrule
G=0, D=5 & 1.16 / 95\% & 2.99 / 95\% & 2.05 / 95\% & 2.90 / 95\% & 2.61 / 95\% & 2.64 / 95\% & 2.47 / 94\% \\
G=5, D=5 & 1.06 / 95\% & 2.22 / 95\% & 1.76 / 95\% & 2.44 / 95\% & 2.01 / 95\% & 1.90 / 95\% & 1.85 / 94\% \\
G=15, D=5 & 1.01 / 95\% & 1.69 / 95\% & 1.47 / 95\% & 1.99 / 95\% & 1.63 / 95\% & 1.55 / 95\% & 1.52 / 94\% \\
G=30, D=5 & 0.90 / 95\% & 1.51 / 95\% & 1.29 / 95\% & 1.63 / 95\% & 1.38 / 95\% & 1.27 / 95\% & 1.33 / 94\% \\
G=60, D=5 & 0.80 / 95\% & 1.32 / 95\% & 1.15 / 95\% & 1.38 / 95\% & 1.12 / 95\% & 1.11 / 95\% & 1.10 / 94\% \\
G=0, D=10 & 1.07 / 90\% & 1.73 / 93\% & 1.50 / 94\% & 1.70 / 94\% & 1.60 / 93\% & 1.53 / 93\% & 1.58 / 91\% \\
G=0, D=15 & 0.89 / 84\% & 1.13 / 93\% & 1.07 / 89\% & 1.13 / 90\% & 1.20 / 90\% & 1.09 / 90\% & 1.09 / 87\% \\
\bottomrule
\end{tabular}}
\end{table*}


FAPD counts alarm episodes, so it depends on what ends an episode. The default
policy ends an alarm as soon as the score dips below threshold, so a single
5-minute dip produces a second sounding. Table~\ref{tab:fapd_latching} re-scores
every arm under two alternatives, holding each seed's calibration-fitted
threshold fixed: latching across gaps of up to $G$ minutes, and a minimum
duration of $D$ minutes before an alarm sounds. Absolute FAPD is highly
policy-dependent: fusion falls from $\NumLatchBaseFusionFAPD$ to
$\NumLatchGSixtyFusionFAPD$ false alarms per day under 60-minute latching, so
FAPD values are comparable only within a policy. Persistence has the lowest FAPD
under every policy tested, at the cost of the shortest onset lead. Fusion's lower
FAPD than LightGBM under latching does not come free: under a 15-minute minimum
duration, fusion detects $\NumLatchDFifteenFusionDet\%$ of onsets against
LightGBM's $\NumLatchDFifteenLGBMDet\%$ (difference $\NumLatchDFifteenDetDiffLGBM$
points, 95\% subject-bootstrap CI $[\NumLatchDFifteenDetDiffLGBMLo,
\NumLatchDFifteenDetDiffLGBMHi]$).

\subsection{Ablation F: Is the Endpoint Baseline Handicapped?}
\label{sec:ablation_f}

\begin{table*}[t]
\centering
\small
\setlength{\tabcolsep}{4pt}
\caption{Ablation F: is the endpoint baseline handicapped by its risk mapping? (PH = 30 min). Classification metrics only. Persistence, Ridge-AR and ARIMA have no classification head and score hypoglycemia through $\sigma((70 - \hat{y})/\tau)$, a strictly monotone map of their forecast; the label is the window \emph{minimum}, which an endpoint forecast does not target. The nadir arms are the identical Ridge-AR fitted to $\min g$ over $(t, t+\mathrm{PH}]$ instead of $g(t+\mathrm{PH})$ and mapped the same way, so the difference between the two rows is the target alone. LightGBM and every neural arm fit a classifier on the nadir label directly and are shown for reference, not as part of the contrast. All ablation arms share one budget (20 epochs, fit split subsampled $2\times$), lighter than Table~\ref{tab:sota_ph30}'s (30 epochs, full fit split); rows here compare to each other, not to Table~\ref{tab:sota_ph30}.}
\label{tab:ablation_f_nadir_target_ph30}
\resizebox{\ifdim\width>\linewidth\linewidth\else\width\fi}{!}{%
\begin{tabular}{lrrrrr}
\toprule
Arm & PR-AUC & ROC-AUC & Se@5\% & FAPD@5\% & Onset lead (min) \\
\midrule
Ridge-AR, level target & 0.829 & 0.9900 & 0.964 & 2.99 & 26.4 \\
Ridge-AR, nadir target & 0.861 & 0.9900 & 0.962 & 2.25 & 26.0 \\
LightGBM (native nadir classifier) & 0.841 & \textbf{0.9916} & \textbf{0.968} & 2.90 & \textbf{27.1} \\
CausalCGM-MT (native nadir head) & \textbf{0.865} & 0.9912 & 0.959 & 2.87 & 26.4 \\
\bottomrule
\end{tabular}}
\end{table*}

\begin{table*}[t]
\centering
\small
\setlength{\tabcolsep}{4pt}
\caption{Ablation F: is the endpoint baseline handicapped by its risk mapping? (PH = 60 min). Classification metrics only. Persistence, Ridge-AR and ARIMA have no classification head and score hypoglycemia through $\sigma((70 - \hat{y})/\tau)$, a strictly monotone map of their forecast; the label is the window \emph{minimum}, which an endpoint forecast does not target. The nadir arms are the identical Ridge-AR fitted to $\min g$ over $(t, t+\mathrm{PH}]$ instead of $g(t+\mathrm{PH})$ and mapped the same way, so the difference between the two rows is the target alone. LightGBM and every neural arm fit a classifier on the nadir label directly and are shown for reference, not as part of the contrast. All ablation arms share one budget (20 epochs, fit split subsampled $2\times$), lighter than Table~\ref{tab:sota_ph30}'s (30 epochs, full fit split); rows here compare to each other, not to Table~\ref{tab:sota_ph30}.}
\label{tab:ablation_f_nadir_target_ph60}
\resizebox{\ifdim\width>\linewidth\linewidth\else\width\fi}{!}{%
\begin{tabular}{lrrrrr}
\toprule
Arm & PR-AUC & ROC-AUC & Se@5\% & FAPD@5\% & Onset lead (min) \\
\midrule
Ridge-AR, level target & 0.702 & 0.9650 & 0.864 & 3.72 & 45.1 \\
Ridge-AR, nadir target & 0.757 & 0.9640 & 0.840 & 2.68 & 42.0 \\
LightGBM (native nadir classifier) & 0.755 & \textbf{0.9718} & \textbf{0.867} & 3.76 & \textbf{46.0} \\
CausalCGM-MT (native nadir head) & \textbf{0.772} & 0.9714 & 0.849 & 3.37 & 43.5 \\
\bottomrule
\end{tabular}}
\end{table*}


Section~\ref{sec:baselines} flagged an asymmetry in the benchmark table:
persistence, Ridge-AR and ARIMA are scored on a nadir label through a monotone
map of an endpoint forecast. This study measures what that costs them, by
refitting Ridge-AR to $\min g$ over $(t, t{+}\PH]$ instead of $g(t{+}\PH)$. The
design matrix, the ridge penalty, the residual parameterisation and the risk
mapping are all unchanged; only the regression target moves, so the difference
between the two rows is attributable to the target and nothing else.

\textbf{The handicap is real and we were understating these baselines.}
Targeting the nadir raises Ridge-AR's PR-AUC by $\NumAblFNadirGainHThirty$ at
$\PH = 30$ and $\NumAblFNadirGainHSixty$ at $\PH = 60$, and lowers its
false-alarm rate from $\NumAblFRidgeLevelFAPDHThirty$ to
$\NumAblFRidgeNadirFAPDHThirty$ per day. ROC-AUC barely moves
($\NumAblFRidgeLevelROCAUCHThirty \to \NumAblFRidgeNadirROCAUCHThirty$), which
locates the gain precisely: the endpoint forecast ranks the bulk of the
distribution just as well, and what it was losing was precision on the rare
positive class. More specifically, those are the transient nadirs that recover
before the horizon, which are the events the label is built around and the
endpoint cannot see. At
$\PH = 30$ the nadir-targeted Ridge-AR ($\NumAblFRidgeNadirPRAUCHThirty$) also
overtakes LightGBM's native classifier ($\NumAblFLGBMPRAUCHThirty$), so this is
not a small correction to an also-ran.

\textbf{With the baseline corrected, the network's forewarning advantage is
not established.} Against nadir-targeted Ridge-AR, CausalCGM-MT is ahead by
$\NumAblFMTLeadOverNadirHThirty$ PR-AUC at $\PH = 30$ and
$\NumAblFMTLeadOverNadirHSixty$ at $\PH = 60$ in this single-seed study. Those
gaps are compared here with the seed-to-seed standard deviation of each neural
arm's PR-AUC in the multi-seed tables. At $\PH = 30$ these are
$\NumMSFusionPRAUCHThirtySD$ (fusion), $\NumMSLSTMPRAUCHThirtySD$ (LSTM) and
$\NumMSCGMOnlyPRAUCHThirtySD$ (CGM-only); at $\PH = 60$ they are
$\NumMSFusionPRAUCHSixtySD$, $\NumMSLSTMPRAUCHSixtySD$ and
$\NumMSCGMOnlyPRAUCHSixtySD$. The $\PH = 30$ gap is below all three. The
$\PH = 60$ gap is under two standard deviations of the fusion and LSTM arms
but several times that of the CGM-only arm, and it comes from one seed, so we
do not claim a lead over a well-posed classical baseline. A claim
that rests on a baseline being asked the wrong question is not a claim about the
model.

It follows that the classification columns for persistence, Ridge-AR and ARIMA
in the benchmark tables are lower bounds on what those methods achieve when
targeted at the labelled quantity, and should not be cited as the ceiling for
classical forecasting.

\subsection{What the Sigmoid Mapping Measures}
\label{sec:appendix_sigmoid}

The mapping asks those three arms an uneven question. The label is $\min g$ over
$(t, t{+}\PH]$; a level forecast at $t{+}\PH$ is blind to a nadir at $t{+}20$
that recovers by $t{+}30$, and no monotone map of that forecast can recover what
the forecast did not encode. Their classification columns therefore measure two
things at once: how well the arm forecasts, and how much nadir information a
$t{+}\PH$ endpoint carries. Section~\ref{sec:ablation_f} separates them by
fitting the same arm to the nadir directly.
